\chapter{The complete positive-scale C1 reciprocal terminal}
\label{component:p98f24864b071}
\noindent\textbf{Component }\path{p98f24864b071}\par
{\small\noindent Source: \path{NEW/global/c1_reciprocal_full_parent_REDERIVED.tex}\par}
{\small\noindent SHA256: \path{98f24864b0712df08a74a0cf16c9b7e8412282a9a8aa6494e7d2f657042b38f3}\par}
{\small\noindent Inventory use: CR1-CR39 full C1 reciprocal all-tests theorem, with AU19 proof-use audit | CR1-CR39 complete source entry, all local states, full normal restoration and projective outer join\par}
\medskip
\section{The complete positive-scale C1 reciprocal terminal}
\label{p98f24864b071:sec:c1-reciprocal-rederived}

\paragraph{Theorem and exact scope.}
Let the canonical C1 source be physical, with an actual finite
noncommuting zero contact, and normalize $d=1$. Suppose every source
coefficient is bounded, $\lambda=n=0$, $C,A,\alpha$ are positive
bounded units, $m>0$ and $D\geq0$ are infinitesimal, and
$\beta\geq0$ is arbitrary bounded. Retain the original canonical
source gaps and intersections. Then the entire parent has a finite
sum-of-squares representation in $K_c\otimes_{\mathbb R}K_t$ for
every independent projective test and every real or complex state.
The assertion includes literal $D=0$ and literal $\beta=0$ when
permitted by the actual contact; they cannot both vanish.

We use $O(1)$ and $o(1)$ for bounded and infinitesimal elements of
one of the two fields, and $a\asymp b$ for positive source quantities
whose ratio is a bounded unit. A bounded tensor Gram with a strictly
positive real residue is a sum of squares by the actual separated
strict-Gram lemma. Its rows need not be bounded. No order on the tensor
product, mixed inverse, or cancellation of a determinant multiplier is
used below.

\subsection{Source entry before any state-ratio division}

Write the actual contact angles as $u_c=k>0$, $b_c=1/v_c=p>0$,
and put $U=k^2$, $P=p^2$. The actual C1 source balances SP2--SP3
give $U,P$ positive infinitesimals and
\[
 D/P^2=O(1),\qquad \beta U/P^2=O(1),\qquad
 m/U^2\asymp1+\beta/P.
 \tag{CR1}
\]
Indeed the two stationarity equations are
\[
\begin{aligned}
 U^2(a_6+\beta a_7/P)&=m(a_1+a_2D/P),\\
 a_3&=P^{-2}(mDa_2/U+\beta ma_4+Da_5+\beta a_7U),
\end{aligned}
\]
where every $a_i$ is a positive bounded unit. Positivity bounds
each term of the second equation; then $D/P=o(1)$ and the
parentheses in the first give CR1. In particular
\[
 m/U=o(1),\qquad D/P=o(1),\qquad \beta U/P=o(1).
 \tag{CR2}
\]
The original source partition proves $U,P$ infinitesimal in exactly
the present $m,D$ weak branch. Only those displayed source balances,
not its absent terminal references, are an input here.

At a physical noncommuting contact the signed core is negative, and
the exact core Schur identity SE9 gives $W_cZ_c>0$. Normalize $W_c=1$
and write $Z_c=q>0$. Define
\[
 Q=q^2,\quad L=PQ,\qquad
 r=\frac{k}{pq},\quad y=\frac{Y_c}{pq},\quad
 x=\frac{X_c}{pq^2}.
 \tag{CR3}
\]
We now prove these divisions are bounded wherever they are used.
In the reciprocal quarter-parent CR9 below the hard $X,Y$ matrix at
the source has residue
$\operatorname{diag}(\bar\beta+1/2,H_0)$, where
$H=C+A+\alpha+1/2$ is a positive unit. Eliminate this two-by-two
block inside the source field. Its Schur block on $W,Z$ satisfies
\[
 \frac{S_{WW}}{U}\sim a:=C-C^2/H>0,\qquad
 -\frac{S_{WZ}}{kp}\sim b_*:=1/2+C\alpha/H>0.
 \tag{CR4}
\]
Here and below a displayed positive expression in an asymptotic
formula is evaluated in residue. To check CR4 without a shape
assumption, the original hard couplings are
\[
\begin{aligned}
 P_{XW}&=p\{-m+(D+1/2)U\},&
 P_{YW}&=k\{m-C(1+P)-(D+1/2)P\},\\
 P_{XZ}&=k\{\beta-D+P(m+\beta+1/2)\},&
 P_{YZ}&=p\{D-\alpha-U(m+\alpha+1/2)\},\\
 P_{WZ}&=-(m+D+1/2)kp.
\end{aligned}
\]
CR2 shows that the $X$ corrections in CR4 vanish in residue and
the $Y$ corrections are respectively $C^2/H$ and $C\alpha/H$.
The kernel equation gives $q=S_{WW}/(-S_{WZ})$, so $r$ is a
positive bounded unit.

The scaled $X,Y$ kernel equations in the unknowns $x,y$ have
bounded entries and a diagonal positive residue. For example their
$X$ equation is exactly
\[
\begin{split}
 0={}&\{(\beta+1/2)(1+P)+Pm+(D+1/2)r^2PQ\}x\\
 &+r\{\beta-D+P[\beta+1/2-m(y-1)+(D+1/2)(r-y)]\}-m/Q.
\end{split}\tag{CR5}
\]
The quantity $m/Q=(m/U)r^2P$ is bounded before this equation is
used. The $y$ coefficient in CR5 is infinitesimal; the $x$
coefficient in the scaled $Y$ equation is $O(PQ)$ and its $y$
pivot is $H+o(1)$. Thus $x,y$ are bounded. Taking the $Y$ and
$Z$ kernel equations in residue, using CR1--CR2, yields
\[
 Hy\sim Cr+\alpha,\qquad C(r-y)\sim1/2,
 \qquad \alpha(1-y)\sim r/2.
 \tag{CR6}
\]
Consequently $y,r-y,1-y$ are positive bounded units. Equivalently
their real residues satisfy
\[
 C_0=\frac1{2(r_0-y_0)},\quad
 A_0=\frac{1+r_0-y_0}{2y_0},\quad
 \alpha_0=\frac{r_0}{2(1-y_0)}.
 \tag{CR7}
\]
For the last equation in CR6, divide the original $Z$ kernel
equation by $Pq$. Its extra $X$ contribution is
$O(x\beta Q+xPQ)$; CR1 and the already proved unit $r$ give
$\beta Q=O(P)$, while $PQ=U/r^2$ is infinitesimal. All the
omitted terms therefore vanish in residue. This proves CR6 in the
unbounded-$q$ case as well.

In CR5, $D=o(P)$ and $m/Q=o(P)$. Its numerator for $-x$ is
the sum of $r\beta$ and a positive-unit multiple of $P$, up to
$o(P)$. Hence the following bounds hold on the whole source domain:
\[
\begin{gathered}
 L=PQ=o(1),\qquad \beta Q/P=O(1),\qquad
 -x\asymp P+\beta,\qquad r+x\text{ is a positive unit},\\
 m\asymp PQ^2(P+\beta),\qquad
 m/(PQ)=O(PQ+\beta Q)=o(1),\qquad xQ=o(1).
\end{gathered}\tag{CR8}
\]
The statement about $r+x$ follows directly from CR5:
$x/r=-\beta/(\beta+1/2)+o(1)$.

\subsection{The entire scalar and all normal coefficients}

Use $b=1/v$, $D_u=1+u^2$, $D_b=1+b^2$. The full reciprocal
quarter-parent is
\[
\begin{aligned}
 F={}&m(W-bX+uY-ubZ)^2+CD_b(uW-Y)^2+AD_uD_bY^2\\
 &+(D+1/2)(ubW+uX-bY-Z)^2
 +(\alpha+1/2)D_u(Y-bZ)^2\\
 &+(\beta+1/2)D_b(X+uZ)^2-D_uD_bZ^2/2.
\end{aligned}\tag{CR9}
\]
All further local identities are expansions of this full expression.
In the chart $W=1$, $Z=z$ put
\[
 u_g=rbz,\qquad X_g=xbz^2,\qquad Y_g=ybz,
 \qquad V=b^2,\quad T=z^2.
 \tag{CR10}
\]
Its entire scalar restriction is
\[
 \Phi=m+DT+hVT+\mathcal C V^2T+EVT^2+
 \mathcal F V^2T^2+GV^3T^2+JV^2T^3,
 \tag{CR11}
\]
where
\[
\begin{aligned}
 \mathcal C&=Ay^2+(C+D+1/2)(r-y)^2,\\
 E&=\beta(r+x)^2+x^2/2-2Drx,\qquad
 G=Ar^2y^2,\qquad J=(D+1/2)r^2x^2,\\
 \mathcal F&=Ar^2y^2+(\alpha+1/2)r^2(y-1)^2
       +\beta(r+x)^2+m[r(y-1)-x]^2\\
 &\hspace{12mm}+(2D+1)rx(r-y)+rx+x^2/2,\\
 h&=Ay^2+C(r-y)^2+\alpha(y-1)^2+y^2/2-r
          -2D(r-y)+2m[r(y-1)-x].
\end{aligned}\tag{CR12}
\]
In particular $\mathcal C,G$ are positive units,
$E\asymp\beta+P^2>0$, $J\geq0$, and all these coefficients
are bounded. If $\beta$ is infinitesimal, $\mathcal F$ is a
positive unit, with residue $r_0^2(1+r_0-y_0)/2$.
For unit $\beta$, $\mathcal F$ can have either sign; it is not
discarded or assumed positive.

Actual contact and its two angular/state jets give
$\Phi(P,Q)=\Phi_V(P,Q)=\Phi_T(P,Q)=0$. Therefore
\[
\begin{aligned}
 D&=P^2\mathcal C-PQE+P^3QG-P^2Q^2J,\\
 m&=PQ^2E+P^2Q^2\mathcal F+P^3Q^2G+2P^2Q^3J,\\
 h&=-2P\mathcal C-QE-2PQ\mathcal F-3P^2QG-2PQ^2J.
\end{aligned}\tag{CR13}
\]
Set
\[
 \delta=D/P^2\geq0,\qquad \theta=EQ/P>0.
 \quad \delta,\theta=O(1),\qquad
 \delta+\theta=\mathcal C+GL-JQ^2=\mathcal C+o(1),
 \tag{CR14}
\]
where $L=PQ$. The bounds follow either from CR1 and CR8 or from CR13;
$JQ^2=O((xQ)^2)=o(1)$. At least one of $\delta,\theta$
is a unit. Define
\[
 R=z(V-P),\quad K=VT-PQ,\quad
 M=V^2T+(P/Q)VT^2+P^2Q-3PVT.
\]
The exact complete balanced scalar is
\[
\begin{aligned}
 \Phi={}&\delta R^2+\theta M+\mathcal F K^2
       +G\{PK^2+VzRK\}\\
       &+J\{(QR-zK)^2+2QK^2\}.
\end{aligned}\tag{CR15}
\]
The entire $J$ line is retained as positive squares.

Here are full, rather than selected, normal formulas. Let $L_u,L_X,L_Y$
be half the first derivatives of CR9 at CR10. Then
\[
 L_u=bz A_u(V,T),\quad L_Y=bz A_y(V,T),\quad
 L_X=b\{-m+TB_X(V,T)\},
 \tag{CR16}
\]
where
\[
\begin{aligned}
 A_u&=a_{00}+a_{10}V+a_{01}T+a_{11}VT+a_{21}V^2T+a_{12}VT^2,\\
 A_y&=e_{00}+e_{10}V+e_{11}VT+e_{21}V^2T,\\
 B_X&=c_0+c_1V+c_2VT,\\
 c_0&=-Dr+\beta(r+x)+x/2,\\
 c_1&=(D+1/2)r(r-y)+(\beta+1/2)(r+x)
                        +m[x-r(y-1)],\\
 c_2&=(D+1/2)r^2x.
\end{aligned}\tag{CR17}
\]
Every displayed coefficient is bounded. The only bounds beyond this
needed below are
\[
 a_{01}=\beta(r+x)-Dx=O(E),\quad
 a_{12}=(D+1/2)rx^2,\quad a_{12}Q=O(1),
 \qquad c_2=O(P+E).
 \tag{CR18}
\]
They follow from CR8 and CR12. The other coefficients are given
unambiguously by expanding CR9; the checker compares the entire
polynomials, not a support truncation.

State and angle stationarity give
$A_u(P,Q)=A_y(P,Q)=0$ and $B_X(P,Q)=m/Q$.
Put $E_T=bz(T-Q)$. Finite monic division therefore gives exactly
\[
\begin{aligned}
 L_i&=A_{i,V}(V,T)bR+B_i(T)E_T,\
 B_u(T)&=a_{01}+P\{a_{11}+Pa_{21}+a_{12}(T+Q)\},\\
 B_y(T)&=P(e_{11}+Pe_{21}),\\
 L_X&=b\{[m/(PQ)+c_2T]K+z[c_1-m/(PQ)]R\}.
\end{aligned}\tag{CR19}
\]
The first line is for $i=u,y$ and
$A_{i,V}=[A_i(V,T)-A_i(P,T)]/(V-P)$, an exact polynomial.
In particular, using $P(T-Q)=K-zR$, each of $L_u,L_Y,L_X$
is a bounded linear combination of
\[
 bR,\qquad bzK,\qquad b[m/(PQ)+c_2T]K,
 \qquad E E_T.
 \tag{CR20}
\]
Every occurrence of $K$ outside the last product has an
infinitesimal coefficient. CR19 remains literal at $b=0$ or $z=0$.

\subsection{The cubic conductors and the bounded source-ratio chart}

All radicals in this subsection are products of radicals in the
separate fields. Set
\[
 \mathsf A=(V^2T)^{1/3},\quad
 \mathsf B=((P/Q)VT^2)^{1/3},\quad
 \mathsf C=(P^2Q)^{1/3},\qquad \rho=(P/Q)^{1/3}.
\]
Because $\mathsf A\mathsf B\mathsf C=PVT$,
\[
 M=\tfrac12(\mathsf A+\mathsf B+\mathsf C)
 [ (\mathsf A-\mathsf B)^2+(\mathsf A-\mathsf C)^2
                         +(\mathsf B-\mathsf C)^2].
 \tag{CR21}
\]
Let $\mathcal Z$ be its nine roots multiplied by $\sqrt\theta$,
so that $\theta M=\frac12\sum_{j=1}^9\mathcal Z_j^2$.
For nonnegative testing $b,z$, the following exact identities hold:
\[
\begin{aligned}
 bR={}&b\sqrt{\mathsf A}(\mathsf A-\mathsf B)
       +\rho b^{1/3}z^{2/3}\sqrt{\mathsf A}(\mathsf A-\mathsf C),\\
 E E_T={}&\sqrt E\,z\,[\sqrt\theta\sqrt{\mathsf B}
                                      (\mathsf B-\mathsf A)]\\
 &+\theta^{1/3}E^{1/6}b^{2/3}z^{1/3}
       [\sqrt\theta\sqrt{\mathsf B}(\mathsf B-\mathsf C)].
\end{aligned}\tag{CR22}
\]
Negative signs are handled by the exact testing gauges; absolute
values give the same monomial identities if desired. If $\theta$
is a unit, the first line has infinitesimal coefficients on
$\mathcal Z$: $\rho=(E/\theta)^{1/3}$ is bounded and $b$ is
infinitesimal. If $\delta$ is a unit, $bR$ instead has the
infinitesimal coefficient $b$ on the strict row $R$. The second
line has infinitesimal coefficients whenever $E$ is infinitesimal
or $z$ is infinitesimal. Its displayed coefficients use no
comparison of a source rate with a testing rate.

Suppose first that actual $q$ is bounded and $\beta$ is
infinitesimal. Use $K$ and $\mathcal Z$, with $R$ added when
$\delta$ is a unit. If only $\theta$ is a unit, retain
$\delta R^2$ as a separate positive square. In CR15 the $K$
pivot $\mathcal F$ is a positive unit. The only signed extra
term is
$GVzRK=Gbz(bR)K$, an infinitesimal Gram correction by CR22.
The term $GPK^2$ is nonnegative. Thus CR15 has a strict residue
Gram on the stated roots. CR19--CR22 give infinitesimal lifts of
all three complete first normals to these same roots.

If $\beta$ is a unit, CR8 forces $Q$ infinitesimal and $E$ a
unit. The only chart still requiring CR10 will be $z$ infinitesimal.
The second line of CR22 still has infinitesimal coefficients.
When $\delta$ is a unit, use the exact identity
\[
 2QM+TR^2-K^2=(K-zR)^2+2QR^2.
 \tag{CR23}
\]
Choose a positive real $a$ with $\overline{\mathcal F}+a>1$.
Subtracting $a$ times the corresponding zero identity changes the
first three terms of CR15 into
\[
 (\delta-aT+2aQ)R^2+(1-2aP/E)\theta M
       +(\mathcal F+a)K^2+a(K-zR)^2.
 \tag{CR24}
\]
The first three pivots have strictly positive residues; both costs
$T$ and $P/E$ are infinitesimal. This handles the signed
$\mathcal F$ exactly.

When instead $\theta$ is a unit, $\rho$ is also a unit and
\[
 K=\rho^{-1}\{\sqrt{\mathsf A}\,
       [\sqrt{\mathsf A}(\mathsf B-\mathsf C)]
       +\sqrt{\mathsf C}\,
       [\sqrt{\mathsf C}(\mathsf A-\mathsf C)]\}.
 \tag{CR25}
\]
Both coefficients on $\mathcal Z$ are infinitesimal. Thus the
entire signed $\mathcal F K^2$ is an infinitesimal correction on
the strict circuit Gram. The $G$ cross term and every first normal
are controlled by CR22 and CR25; $\delta R^2$ and the $J$ line
remain positive squares. No inverse of $\delta$ is used.

\subsection{An unbounded actual state ratio in the same testing chart}

Now let actual $q$ be unbounded. CR8 and CR12 imply
\[
 \beta/P=O(1/Q)=o(1),\quad -x\asymp P,\quad
 m\asymp L^2,\quad E/P=o(1),\quad JQ^2=O(L^2)=o(1).
 \tag{CR26}
\]
Here $\beta$ is infinitesimal and $\mathcal F$ is a positive
unit. In the testing chart $W=1$, $z$ bounded, the exact relation
\[
 \theta M=\theta R^2+(E/P)(K^2-zRK)
 \tag{CR27}
\]
makes CR15 a strict two-row Gram on $R,K$, plus its whole positive
$J$ line. Its residue is
$\operatorname{diag}(\overline{\delta+\theta},\overline{\mathcal F})$.
Every correction coefficient is infinitesimal: $E/P$, $P$ and
$Vz$ are infinitesimal. In CR19, $B_u/P$ and $B_y/P$ are bounded,
because $a_{01}/P=O(\beta/P+Dx/P)$ and $a_{12}Q$ are bounded.
Using $P E_T=bzK-bz^2R$ shows that all three full normals are
infinitesimal linear combinations of $R,K$.

\subsection{The complementary testing state chart}

Use $Z=1$, $W=w$ infinitesimal, and write $a=u/r$. The graph is
\[
 b=aw,\qquad X=xa,\qquad Y=yaw.
 \tag{CR28}
\]
The entire scalar is
\[
 \Phi'=D+mw^2+Ea^2+ha^2w^2+\mathcal C a^4w^4
       +\mathcal F a^4w^2+Ga^6w^4+Ja^4.
 \tag{CR29}
\]
First assume actual $q$ bounded and $\beta$ infinitesimal. If
$\delta$ is a unit, use the three strict rows
$P,a^2w,\sqrt E\,a$. The sole signed cross in CR29 is
$h a^2w^2=(h/P)w\,P(a^2w)$, and $h/P$ is bounded by CR13.
Every other term is positive or has the positive-unit $\mathcal F$
pivot. If $\theta$ is a unit, $P/E=Q/\theta$ is bounded;
use instead $\sqrt E\,a,a^2w$ and retain $D,mw^2$ whole.
The same signed term is the infinitesimal correction
$(h/E)w^2(\sqrt E\,a)^2$.

Here are the complete complementary first normals, which also prove
the bounds for unbounded actual $q$. Define
\[
 R'=a^2w^2-P,\qquad K'=w(a^2-L).
 \tag{CR30}
\]
Half the $b,Y$ derivatives of the original full CR9 at CR28 have
the respective forms
\[
 L_i'=aw(d_{i0}+d_{i1}a^2+d_{i2}a^2w^2+d_{i3}a^4w^2),
 \qquad i=b,Y,
 \tag{CR31}
\]
with bounded coefficients. At the actual contact
$a^2=L,w^2=1/Q$, they satisfy
$d_{i0}+Ld_{i1}+Pd_{i2}+PLd_{i3}=0$.
Consequently the entire, exact divisions are
\[
 L_i'=(d_{i1}+Pd_{i3})aK'
                   +(d_{i2}+a^2d_{i3})awR'.
 \tag{CR32}
\]
The full $X$ derivative is
\[
 L_X'=a(c_0+c_2a^2-mw^2+c_1a^2w^2),\qquad
 c_0+c_2L-m/Q+c_1P=0.
 \tag{CR33}
\]
For bounded $Q$, $c_0/P$ and $m/P$ are bounded, and
$c_2\in P B_c+E B_c$ by CR18. The products $aK'$ and
$awR'$ are infinitesimal linear combinations of $P,a^2w$.
Expanding CR33 gives infinitesimal coefficients on
$P,a^2w,\sqrt E\,a$. If $\theta$ is a unit, use $P/E=O(1)$:
$c_0/\sqrt E=o(1)$, $c_2=O(E)$ and $m/E=O(1)$, giving the
same assertion on the two rows $\sqrt E\,a,a^2w$.
These are pure-source bounded quotients; no testing $w$ is inverted.

For unbounded actual $q$, CR29 has the exact alternative form
\[
\begin{aligned}
 \Phi'={}&\delta R'^2+\theta M'+\mathcal F K'^2
       +G\{PK'^2+a^2wR'K'\}\\
       &+J\{a^4(Qw^2-1)^2+2QK'^2\},\\
 M'={}&a^4w^4+(P/Q)a^2+P^2Qw^2-3Pa^2w^2.
\end{aligned}\tag{CR34}
\]
Define nonnegative elementary radicals
\[
 \mathsf A'=|aw|^{4/3},\quad
 \mathsf B'=(P/Q)^{1/3}|a|^{2/3},\quad
 \mathsf C'=(P^2Q)^{1/3}|w|^{2/3}.
\]
They give the same nine-square cubic identity for $M'$. The exact
conductor needed for the complete normals is
\[
 |a|R'=|a|\sqrt{\mathsf A'}(\mathsf A'-\mathsf C')
          +\sqrt L\sqrt{\mathsf B'}(\mathsf A'-\mathsf B').
 \tag{CR35}
\]
Its coefficients are infinitesimal. If $\delta$ is a unit, use
$R',K'$ and retain the circuit as positive squares. If $\theta$
is a unit, use its scaled nine roots and $K'$, retaining
$\delta R'^2$ whole. In either case CR34 is strict: its sole
signed cross is $Gaw(aR')K'$, controlled by CR35, and its last
line is retained as squares.

CR26 implies $c_2/P=O(1)$ and
$c_0/P=m/(PQ)-Qc_2-c_1=O(1)$. The full exceptional normal
has the exact division
\[
\begin{aligned}
 L_X'={}&-[c_0/P+(c_2/P)a^2]\,aR'\\
 &+[(c_2/P)(a^2+L)+c_0/P+c_1]\,awK'.
\end{aligned}\tag{CR36}
\]
Together with CR32, this uses only bounded multiples of
$aR',awR',aK'$. Each has infinitesimal coefficients on the strict
soft rows by CR35. This also proves the result at literal $w=0$.

\subsection{All higher normals and the remaining unit-beta states}

The preceding scalar and first-normal formulas certify the entire
parent. Here are the required hard blocks and the justification for
every higher term. Normalize the testing state so its coordinates
are bounded with at least one unit. At the reciprocal angle corner
$u,b=o(1)$ the state residue hard block is
$(\bar\beta+1/2)X^2+H_0Y^2$. If $X$ or $Y$ is a unit, the
entire parent has positive real residue and the strict-element lemma
applies immediately. Otherwise $X,Y$ are infinitesimal and one of
$W,Z$ is a unit.

The exact sign gauges make $u,b\geq0$. If $WZ\leq0$, the signed
core is a nonnegative pure-testing scalar by the full core Schur
identity SE9 and its literal homogeneous endpoints, so its source
multiple and all event packets are separated squares. Otherwise
choose either $W=1$, $z=Z$ bounded, or $Z=1$, $w=W$
infinitesimal. Both have $X,Y$ infinitesimal.

In the $W=1$ chart, set
$N_u=u-rbz$, $N_X=X-xbz^2$, $N_Y=Y-ybz$.
They are infinitesimal. The complete expansion of CR9 is its scalar
CR11, twice the three first normals CR16, and a polynomial quadratic
Gram in $N_u,N_X,N_Y$. Terms of normal degree at least three
have infinitesimal Gram entries. For weak $\beta$, or for unit
$\beta$ with $z$ infinitesimal, the hard residue is
\[
 C_0(N_u-N_Y)^2+(A_0+\alpha_0+1/2)N_Y^2
                         +(\beta_0+1/2)N_X^2\succ0.
 \tag{CR37}
\]
All first-normal lifts established above have infinitesimal
coefficients. The full hard-and-soft residue is therefore positive
definite. In the $Z=1$ chart use
$N_b=b-aw$, $N_X=X-xa$, $N_Y=Y-yaw$.
For weak $\beta$ its hard residue is
\[
 N_X^2/2+H_0N_Y^2-2\alpha_0N_bN_Y+\alpha_0N_b^2\succ0,
 \tag{CR38}
\]
since $H_0-\alpha_0=C_0+A_0+1/2>0$.
The same finite expansion handles every higher normal term. Thus
none of CR15, CR27 or CR34 is being used merely as a scalar
restriction.

It remains only unit $\beta$ and a testing state with $Z$ a unit
relative to $W$. Normalize $Z=1$ and $W=\omega$ bounded. Retain
the entire $m,D$ packets. The remainder belongs exactly to the
square of the ideal $(u,b,X,Y)$. Its quadratic real residue is
\[
\begin{aligned}
 &(\beta_0+1/2)\left(X+\frac{\beta_0}{\beta_0+1/2}u\right)^2
      +\frac{\beta_0}{1+2\beta_0}u^2\\
 &\quad+H_0\left(Y-\frac{C_0\omega u+\alpha_0 b}{H_0}\right)^2
      +a_0(\omega u-r_0b)^2,
 \qquad a_0=C_0-C_0^2/H_0>0.
\end{aligned}\tag{CR39}
\]
CR7 is precisely the critical relation needed for this identity.
The expression is positive definite in $(u,b,X,Y)$ for every
real $\omega$, including zero. Its complete bounded polynomial
matrix is consequently strict; every higher degree is an
infinitesimal correction. This completes all local states, including
the unit-beta case where $\mathcal F$ was signed.

\subsection{Every other test and restoration}

The exact outer identities RC18--RC25 from the newly rederived
C2 reciprocal theorem apply here as polynomial identities, with
$n=0$ and the entire extra $\beta$ packet retained. Their strictness
uses only $C,A,\alpha$ positive units and the critical residue
relations CR7; no C2 source membership is used. Explicitly:
\begin{itemize}
\item For bounded original $u,v$, the four rows $(uW,Y,X,Z)$
give RC18. After CR7 its real determinant is
$ (r_0+u^2v^2)D_u^2D_v^2(1+r_0-y_0)^2/
[16y_0(r_0-y_0)(1-y_0)]>0$.
\item For Alice reciprocal and Bob bounded, use the four rows
$(W,Y,vX,sX+Z)$ and the full core RC20. Its real residue RC21
is strict using $C,A>0$.
\item For Bob reciprocal and Alice a bounded unit, use
$(uW-Y,Y,X,bZ,\sqrt D\,Z)$ and the entire $D$ packet as in
RC22--RC23. CR7 makes the three-by-three angular determinant
strict. Literal $D=0$ causes its last row to be zero, which is
allowed in a strict coefficient Gram.
\item For both angles reciprocal, RC24--RC25 give the five rows
$(W,X,Y,bZ,\sqrt D\,sZ)$ and a strictly positive real residue.
\end{itemize}
The only remaining angle chart is exactly $u,b=o(1)$, treated
above. These are whole-parent identities with every omitted event
packet retained, and include all literal axes and infinities.
The exact source and testing sign gauges, positive pure-field
homogeneous normalization factors and real/imaginary state splitting
restore the stated projective tensor-SOS theorem.

\paragraph{Dependencies and verification boundary.}
The mathematical inputs are the actual full physical parent, the
displayed SP2--SP3 balance proof, SE9 and the exact sign gauges,
the separated strict-Gram lemma, and only the explicit outer
polynomial identities RC18--RC25 with their real residue tests.
No unavailable October theorem or claimed historical admission is
a premise. The focused checker verifies the full scalar, all first
normals with explicit contact-jet remainders, the complementary
divisions, the radical and pivot identities, and the complete hard
residues. Source boundedness, sign, the quantified chart cover and
strictness are the analytical arguments above. Other positive-scale
source types, weak-source families and global finite-lift existence
remain separate obligations.
